\chapter{恢复、可靠性和失败路径}

CCB 的可靠性不是“永不失败”，而是失败后能解释、能恢复、能保留 lineage。通信系统尤其如此。provider 会超时，pane 会死亡，reply delivery 可能卡住，mailbox summary 可能损坏，daemon 可能重启。可靠性设计的目标是让这些故障被收敛到可观察、可重试、可清理的状态。

\section{失败分类}

可以把失败分成六类。

\begin{enumerate}
  \item 提交失败：CLI 无法连接 daemon、target 无效、sender 无效、artifact 写入失败。
  \item 排队失败：dispatcher 无法创建 job 或 message-bureau submission。
  \item 启动失败：agent runtime 不可用、pane dead、provider binding 缺失。
  \item 执行失败：API error、transport error、provider process error。
  \item 完成检测失败：timeout、reply 未稳定、provider session 无法读取。
  \item 投递失败：reply record 存在但 caller mailbox head 或 reply delivery 卡住。
\end{enumerate}

每类失败应该落到不同的证据层。把所有失败都显示成 “job failed” 会让用户无法判断下一步该 retry、resubmit、ack、doctor 还是 restart。

\section{automatic retry}

自动 retry 由 finalization 阶段触发。入口在 \src{lib/ccbd/services/dispatcher_runtime/finalization_runtime/message_bureau_retry.py}，策略代码在 \src{finalization_retry_runtime/}。

自动 retry 需要满足：

\begin{itemize}
  \item message retry policy 的 \src{mode} 是 \src{auto}。
  \item completion decision status 是 failed。
  \item failure reason 或 diagnostics error type 属于可重试集合。
  \item 非 nonretryable API failure。
  \item 未达到 max attempts。
  \item runtime failure 场景下 provider 支持 resume，或 policy 允许 runtime retry。
\end{itemize}

默认可重试 reason 包括 \src{api_error} 和 \src{transport_error}。默认 runtime reasons 包括 \src{pane_dead}、\src{pane_unavailable}、\src{runtime_unavailable}、\src{backend_unavailable}。

\section{manual retry}

手动 \cmd{retry} 不直接复制 job。它先解析 attempt，要求 attempt 已终止、不是 completed、且是对应 agent 的最新 attempt。然后它创建新 job，并在 message bureau 中记录 retry attempt。

这种设计防止两个危险操作：

\begin{enumerate}
  \item 对一个已经成功完成的 attempt 重试，导致重复执行。
  \item 对旧 attempt 重试，绕过后续已经产生的新 attempt。
\end{enumerate}

如果旧 job 已经开始 reply，retry 会倾向于用 \src{continue} 作为 request body，并通过 provider options 标记 retry source job。这体现了 provider resume 和 communication lineage 的结合。

\section{comms recover}

通信恢复入口位于 \src{lib/ccbd/services/dispatcher_runtime/comms_recover.py}，handler 位于 \src{lib/ccbd/handlers/comms_recover.py}。它处理的不是普通用户 retry，而是通信链路出现可恢复阻塞时的定向修复。

recoverability 判断覆盖三类场景：

\begin{itemize}
  \item business job 已失败、取消或 incomplete，且可以 retry。
  \item reply delivery job 已失败、取消或 incomplete，且可以 retry。
  \item business job 仍 running，但 runtime 状态、pane state 或健康状态显示它已经 stale。
\end{itemize}

stale running 的判断会看 agent runtime state、runtime health、pane state 和 running hints。可恢复 stale running 会先 cancel old job，且 \src{record_reply=False}，避免为旧卡住任务生成误导性 reply，然后 retry。

\section{reply delivery repair}

reply delivery repair 位于 \src{lib/ccbd/services/dispatcher_runtime/reply_delivery_runtime/repair.py}。它处理 mailbox head 上的 reply delivery 租约问题。

主要场景：

\begin{itemize}
  \item head 有 reply id，但没有 delivery job id。若 lease 指向同一 inbound event，则 reset stale head。
  \item head 指向的 delivery job 不存在。reset stale reply head。
  \item delivery job running，但 persisted execution state 已经有 terminal pending decision。直接 complete。
  \item delivery job running，但 execution service 不认为它 live。用 reply delivery failed decision 完成旧 job，让 head 重新排队。
\end{itemize}

repair 不应吞掉 reply。失败 delivery job terminal 后，\src{resolve_reply_delivery_terminal} 会把 head 重写回 queued，清理 delivery job id 和 progress，让后续 delivery 可以重新 claim。

\section{reply delivery preparation}

\src{prepare_reply_deliveries} 先运行 repair，再遍历 configured agents，寻找 head reply event。对每个 head，它会：

\begin{enumerate}
  \item 读取 reply id。
  \item 检查是否已有 delivery job。
  \item 读取 ReplyRecord。
  \item 计算 project id。
  \item 构造 reply delivery job。
\end{enumerate}

这说明 reply delivery 是普通 dispatcher job 的一种特殊用法。它有 job id、attempt、completion、terminal state，也会进入 trace。

\section{cancel 的恢复语义}

cancel 是一种主动恢复。它要求 job 不是 terminal。cancel 后，dispatcher 会从 queued 和 active state 移除 job，执行 service cancel，写 cancelled completion decision，并记录 message-bureau terminal state。

cancel 的风险是用户可能误以为它会删除所有通信记录。实际上它不会删除 lineage。trace 仍应能看到 cancelled job、attempt state 和相关 events。这是为了审计和后续 retry。
普通 empty cancel 以 consumed completion notice 保留 trace 可见性而不占用
caller mailbox。chain child cancel 会推进 callback continuation；但 stale-job
恢复使用 \src{record_reply=False} 时不会生成 cancelled reply 或 continuation，
因为随后创建的 retry attempt 继续拥有原 message lineage 的终态责任。

\section{mailbox summary degraded}

queue/inbox view 可能遇到 mailbox summary missing 或 error。renderer 会提示 degraded，并建议 \cmd{ccb doctor} 或等待 maintenance refresh。控制层仍会尝试从 inbound event records 重建 detail。

这种降级策略比直接失败更有用，因为 summary 损坏不代表 message/attempt/reply 记录都损坏。用户仍可能通过 trace 找到可恢复路径。

\section{maintenance heartbeat}

maintenance heartbeat 默认 disabled，但配置结构已存在。它的设计目标是周期性评估 agent health，并把 unknown/unhealthy 风险交给 assessor，默认 assessor 是 \src{ccb_self}。

开发者要避免让 heartbeat 成为隐藏的行为改变来源。heartbeat 应报告、调度或显式执行维护动作，而不是悄悄改变用户没有授权的 runtime state。

\section{diagnostics-first 原则}

恢复逻辑应遵循 diagnostics-first 原则：

\begin{enumerate}
  \item 先记录 recoverability 和 block reason。
  \item 如果已经被后续 attempt 修复，返回 noop 和 latest job id。
  \item 如果无法恢复，返回 noop reason。
  \item 如果执行恢复，记录 cancelled old、released event、retried job、next started。
  \item 恢复后再次计算 recoverability。
\end{enumerate}

\src{comms_recover} 的 audit payload 正体现了这个原则。它让自动修复和用户诊断共用同一份事实，而不是只输出“修好了”。

\section{可靠性测试}

恢复路径至少需要覆盖：

\begin{itemize}
  \item failed terminal job retry。
  \item stale running job cancel+retry。
  \item reply delivery failed job requeue。
  \item orphaned reply delivery lease reset。
  \item already retried noop。
  \item completed job 不可 retry。
  \item mailbox summary missing degraded view。
  \item provider runtime unavailable retry policy。
\end{itemize}

当前测试地图中有 \src{test/test_ccbd_comms_recover.py}、\src{test/test_reply_delivery_formatting.py}、\src{test/test_reply_delivery_start_completion.py}、\src{test/test_ccbd_retry_failure_detail.py}、\src{test/test_message_bureau_control_queue.py} 等相关入口。后续应补齐 repair 场景的显式索引。
